%% ma: L10-B02-0021
%% lop: 10
%% chuong: 1
%% bai: 2
%% dang: 10.2.7
%% loai: TLN
%% muc: VD
%% dap_an: "4"
%% nguon: "Ảnh giáo viên gửi 10/10/2026 (ThS. Nguyễn Hoàng Việt, Chương 2 Tập hợp), A-Đề 1, Phần 3 Câu 18 (Câu 22 giống hệt)"
%% kien_thuc: "Bài 2 (tập hợp và các phép toán trên tập hợp)"
%% trang_thai: chinh-thuc
\begin{ex}
Cho tập hợp $A=[4;7]$ và $B=(2a-1;3a+5]$ khác tập rỗng, với $a\in\mathbb R$. Gọi $S$ là tập hợp các giá trị nguyên của $a$ để $A\subset B$. Có bao nhiêu tập con của $S$?
\shortans{4}
\loigiai{$B\ne\varnothing\Leftrightarrow2a-1<3a+5\Leftrightarrow a>-6$. $A\subset B\Leftrightarrow2a-1<4$ và $3a+5\ge7\Leftrightarrow\dfrac23\le a<\dfrac52$. Kết hợp: $\dfrac23\le a<\dfrac52$. Giá trị nguyên: $a\in\{1;2\}$, nên $S=\{1;2\}$ có $2^2=4$ tập con.}
\end{ex}
